% Shared statements (bodies only). Each document wraps them in its own theorem environments and labels.
% ---------------------------------------------------------------------------------------------
% Definitions and assumptions
% ---------------------------------------------------------------------------------------------
\newcommand{\StmtClass}{%
Fix constants $0<\lambda\le\Lambda$, $C_\ast\ge1$, $c_\ast\in(0,1]$. For $\varepsilon\in(0,1]$, a random vector $X$ in $E$
belongs to the class $\mathcal X(\varepsilon)$ if there are a covariance $\Sigma$ with
$\lambda\,\Id\preceq\Sigma\preceq\Lambda\,\Id$, a symmetric $3$-tensor $\cK$ with $|\cK|\le C_\ast\varepsilon$, and an even
polynomial $\widehat r$ of degree at most $6$ with $\widehat r(0)=0$ and coefficients bounded by $C_\ast\varepsilon^2$,
such that, with the \emph{comparison transform}
$\widehat\psie(t)=e^{-\inner t{\Sigma t}/2}\big(1+\tfrac16\cK[(it)^{\otimes3}]+\widehat r(t)\big)$ and
$\chi_X(t)=\E e^{i\inner tX}$, $a=c_\ast/\varepsilon$:
\begin{itemize}
\item[\textup{(E1)}] $\int_{t_0+[-a,a]^N}|\chi_X(t)|\,dt\le C_\ast$ for every $t_0\in E$;
\item[\textup{(E2)}] $\int_{[-a,a]^N}|\chi_X(t)-\widehat\psie(t)|\,dt\le C_\ast\varepsilon^3$;
\item[\textup{(E3)}] $\Prob(|X|>u)\le C_\ast e^{-u^2/C_\ast}$ for all $u\ge0$.
\end{itemize}
The \emph{comparison density} is the signed density $\psie$ with $\int e^{i\inner ty}\psie(y)dy=\widehat\psie(t)$; explicitly
$\psie=\varphi_\Sigma\,(1+\frac16\inner{\cK}{\He_3^\Sigma}+r)$ with $r$ an even polynomial. Constants denoted $C,c$ in
statements about $\mathcal X(\varepsilon)$ depend only on $N$, $\nu$, $\lambda,\Lambda,C_\ast,c_\ast$.}

\newcommand{\StmtAssB}{%
$\zeta,\zeta_1,\zeta_2,\dots$ are i.i.d.\ in $E$ with $|\zeta|\le M$ almost surely, $\E\zeta=0$, and covariance
$\Sigma\succeq\lambda\,\Id$, $\lambda>0$; $K_3=\E\zeta^{\otimes3}$. A \emph{weight sequence} is $w=(w_k)_{k\ge1}$ with
$w_k\ge0$, $\sum_kw_k=1$; we set
\[
\rho_p=\sum_kw_k^p,\quad\ \xi_w=\sum_kw_k\zeta_k,\quad\ X_w=\rho_2^{-1/2}\xi_w,\quad\
\ew=\frac{\max_kw_k}{\sqrt{\rho_2}},\quad\ \Lw=\log(e/\ew).
\]
No condition is imposed on the characteristic function of $\zeta$.}

\newcommand{\StmtAssR}{%
There is a closed cone $V\subset E$ such that $o\in C^\infty(E\setminus V)$,
$|D^ko(y)|\le C_k\dist(y,V)^{-k}$ for all $k\ge0$, and
$\vol\{y\in B(0,1):\dist(y,V)\le r\}\le C_Vr^{c}$ for $r\le1$.}

\newcommand{\StmtAssL}{%
$N\ge3$, and there are $\kappa\in(0,1]$ with $2+\kappa<N$ and a constant $C$ such that the Fourier transform of the
finite signed measure $\lambda_\nu=\Delta(\chi\nu)$ satisfies $|\widehat{\lambda_\nu}(t)|\le C(1+|t|)^{-\kappa}$ for all
$t\in E$. Here $\chi$ is a fixed smooth radial cut-off equal to $1$ on $B(0,1)$ and supported in $B(0,2)$.}

\newcommand{\StmtAssN}{%
$\xi_b=\frac1b\sum_{i\le b}\zeta_i$ with $\zeta_i$ i.i.d.\ in $E$, $\E\zeta=0$, $\E|\zeta|^3<\infty$, covariance
$\Sigma\succ0$, third cumulant tensor $K_3=\E\,\zeta^{\otimes3}$, and Cram\'er's condition
$\limsup_{|t|\to\infty}|\E e^{i\inner t\zeta}|<1$.}

% ---------------------------------------------------------------------------------------------
% Paper II statements
% ---------------------------------------------------------------------------------------------
\newcommand{\StmtCodim}{%
Let Assumption~\ref{ass:R} hold with $c\ge2$. There are $\varepsilon_0,C>0$ such that for every
$X\in\mathcal X(\varepsilon)$ with $\varepsilon\le\varepsilon_0$, writing $\Le=\log(e/\varepsilon)$,
\begin{equation}\label{eq:codimX}
\argmin_{x\in E}\ \E\big[\nu(x+X)-\nu(X)\big]\ \subset\ B\big(\hat x,\ C\varepsilon^{\min(c,3)}\Le^{N/2}\big),\qquad
\hat x=-\tfrac16J_\Sigma^{-1}T_\Sigma[\cK].
\end{equation}
In particular, under Assumption~\ref{ass:B}, for every weight sequence with $\ew\le\varepsilon_0$,
\begin{equation}\label{eq:codimw}
\arg\min\psi_{\xi_w}\ \subset\ B\Big(\frac{\rho_3}{\rho_2}\,\gs\,,\ C\sqrt{\rho_2}\;\ew^{\min(c,3)}\,\Lw^{N/2}\Big),
\end{equation}
so that for the mean of $b$ samples
$\Med_\nu(\xi_b)=-\gs/b+O(b^{-3/2}\log^{N/2}b)$ if $c=2$ and $\Med_\nu(\xi_b)=-\gs/b+O(b^{-2}\log^{N/2}b)$ if $c\ge3$.}

\newcommand{\StmtLaplace}{%
Let Assumption~\ref{ass:L} hold and put $q=N/(N-\kappa)\in(1,\frac32]$. There are $\varepsilon_0,C>0$ such that for every
$X\in\mathcal X(\varepsilon)$ with $\varepsilon\le\varepsilon_0$, with $H(x)=\E[\nu(x+X)-\nu(X)]$,
$\Psi(x)=\int(\nu(x+y)-\nu(y))\psie(y)dy$ and $d_\varepsilon=\varepsilon^{q}\Le^{(N+1)/4}$:
\begin{equation}\label{eq:laplaceX}
\sup_{|x|\le2}|H(x)-\Psi(x)|\le Cd_\varepsilon^2,\qquad
\sup_{|x|\le1}\ \sup_{v\in\partial H(x)}|v-\nabla\Psi(x)|\le Cd_\varepsilon,\qquad
\arg\min H\subset B(\hat x,Cd_\varepsilon),
\end{equation}
with $\hat x=-\frac16J_\Sigma^{-1}T_\Sigma[\cK]$. In particular, under Assumption~\ref{ass:B}, for every weight sequence
with $\ew\le\varepsilon_0$,
\begin{equation}\label{eq:laplacew}
\arg\min\psi_{\xi_w}\ \subset\ B\Big(\frac{\rho_3}{\rho_2}\,\gs\,,\ C\sqrt{\rho_2}\;\ew^{q}\,\Lw^{(N+1)/4}\Big),
\end{equation}
so that $\Med_\nu(\xi_b)=-\gs/b+O\big(b^{-1-\frac{\kappa}{2(N-\kappa)}}\log^{\frac{N+1}4}b\big)$ for the mean of $b$ samples.
The middle estimate of \eqref{eq:laplaceX} holds for \emph{every} subgradient, hence for the expectation of every
measurable selection of $\partial\nu(x+X)$, whatever its values on the singular set.}

\newcommand{\StmtNuclear}{%
Let $\nu=\nnuc\cdot$ on $\R^{K\times d}$ with $(K,d)\ne(1,1)$, $N=Kd$, and let Assumption~\ref{ass:B} hold. Then for every
bounded law with nondegenerate covariance, lattice laws included,
\[
\Med_\nu(\xi_b)=-\frac{\gs}b+O\big(b^{-1-\gamma}\log^{N}b\big),\qquad
\gamma=\begin{cases}1&\text{if }|K-d|\ge2,\\[1pt] \tfrac12&\text{if }|K-d|=1,\\[1pt]
\tfrac1{14}&\text{if }K=d=2,\\[1pt] \tfrac1{4K^2-2}&\text{if }K=d\ge3,\end{cases}
\]
and more generally $\arg\min\psi_{\xi_w}\subset B\big(\frac{\rho_3}{\rho_2}\gs,\,C\,(\max_kw_k)\,\ew^{2\gamma}\Lw^{N}\big)$,
uniformly over weight sequences with $\ew\le\varepsilon_0$. For $K=d=1$ the
statement is false (Theorem~\ref{thm:sign}).}

\newcommand{\StmtSign}{%
Let $\zeta$ be a real random variable with $\E\zeta=0$, variance $\sigma^2>0$, $\E|\zeta|^3<\infty$,
$\kappa_3=\E\zeta^3$, and $\gs=\kappa_3/(6\sigma^2)$; let $S_b=\zeta_1+\dots+\zeta_b$ and
$G_b^\st=\arg\min_G\E|G+S_b/b|$.
\begin{enumerate}
\item[\textup{(a)}] If $\zeta$ is not lattice, $b\,G_b^\st\to\gs$ (every selection).
\item[\textup{(b)}] If $\zeta$ is lattice, supported on $a+h\Z$ with maximal span $h$, let $\Lambda_b=-(ba+h\Z)$. For every
$\delta>0$ and $b\ge b_0(\delta)$: $b\,G^\st_b\subset[\gs-\frac h2-\delta,\ \gs+\frac h2+\delta]$, and if a point
$\ell$ of $\Lambda_b$ satisfies $|\ell-\gs|\le\frac h2-\delta$ then $b\,G^\st_b=\{\ell\}$. Consequently:
\begin{enumerate}
\item[\textup{(i)}] if $a/h$ is irrational, the limit points of $b\,G_b^\st$ fill the interval $[\gs-\frac h2,\gs+\frac h2]$;
\item[\textup{(ii)}] if $a/h=p/q$ in lowest terms, then for $b$ in each residue class $j$ modulo $q$ for which $\gs$ is not
a midpoint of $-(\frac{jph}q+h\Z)$, $b\,G_b^\st$ is eventually equal to the point of that coset nearest to $\gs$; there
is no limit if $q\ge2$, and if $q=1$ and $\gs\notin\frac h2+h\Z$ the limit is $h\operatorname{round}(\gs/h)$;
\item[\textup{(iii)}] if $\zeta\in h\Z$ and $|\kappa_3|<3\sigma^2h$, then $G_b^\st=\{0\}$ for all large $b$: sign descent is
exactly unbiased, although $\gs\ne0$ when $\kappa_3\ne0$.
\end{enumerate}
\end{enumerate}
In the lattice case $b\,G_b^\st\to\gs$ only if $a\in h\Z$ and $\gs\in h\Z$.}

\newcommand{\StmtWor}{%
Let $n\ge3$ and $z_1,\dots,z_n\in E$ with $\sum_iz_i=0$, $|z_i|\le M$, $\Sigma_n=\frac1n\sum_iz_iz_i^\top\succeq\lambda\,\Id$,
$K_{3,n}=\frac1n\sum_iz_i^{\otimes3}$, and let $\gs^{(n)}=-\frac16J_{\Sigma_n}^{-1}T_{\Sigma_n}[K_{3,n}]$. Let
$\xi_{n,b}$ be the mean of a uniformly random subset of size $b\in\{1,\dots,n-1\}$, $p=b/n$ and $v=np(1-p)=b(n-b)/n$. Suppose that $\nu$
satisfies Assumption~\ref{ass:R} with $c\ge2$, and put $\gamma=\frac12(\min(c,3)-1)$, or Assumption~\ref{ass:L}, and put
$\gamma=\frac{\kappa}{2(N-\kappa)}$. Then there are $v_0,C$, depending only on $\nu,M,\lambda$, such that whenever
$v\ge v_0$,
\begin{equation}\label{eq:wor}
\sup_{m\in\Med_\nu(\xi_{n,b})}\Big|\,b\,m+\frac{n-2b}{n-2}\,\gs^{(n)}\Big|\le C\,v^{-\gamma}\,\log^{N}(ev).
\end{equation}
In particular the first-order bias of a batch drawn without replacement is that of a batch drawn with replacement
multiplied by $\frac{n-2b}{n-2}$: it vanishes at $b=n/2$ and \emph{changes sign} beyond, for every bounded population
with nondegenerate covariance, lattice populations included, as soon as $b\to\infty$ and $n-b\to\infty$.}

\newcommand{\StmtMomCor}{%
Let $N_\beta=(1-\beta)\sum_{k\ge0}\beta^k\xi_b^{(k)}$ with independent batch means $\xi_b^{(k)}$ of samples satisfying
Assumption~\ref{ass:B}, and let $\beff=b\,\frac{1+\beta+\beta^2}{1-\beta^2}$, $\theta=\frac{1-\beta^2}b$. If $\nu$ satisfies
Assumption~\ref{ass:R} with $c\ge2$ (put $\gamma=\frac12(\min(c,3)-1)$) or Assumption~\ref{ass:L} (put
$\gamma=\frac\kappa{2(N-\kappa)}$), then
\[
\Med_\nu(N_\beta)=-\frac{\gs}{\beff}+O\Big(\frac{\theta^{\gamma}\log^{N}(e/\theta)}{\beff}\Big)
\qquad\text{uniformly in }b\ge1,\ \beta\in[0,1)\text{ with }\theta\le\varepsilon_0^2 .
\]
In particular the first-order law holds at \emph{fixed} batch size, even $b=1$, as $\beta\to1$.}

% ---------------------------------------------------------------------------------------------
% Paper I statements
% ---------------------------------------------------------------------------------------------
\newcommand{\StmtConv}{%
Let $F:E\to\R$ be differentiable, strongly convex, with Lipschitz gradient; let $\nu$ be any norm on $E$ and
$(\xi_t)_{t\ge1}$ i.i.d.\ with $\E|\xi_t|^2<\infty$. From deterministic $W_1,M_0$ run
\[
M_t=\beta M_{t-1}+(1-\beta)\big(\nabla F(W_t)+\xi_t\big),\qquad W_{t+1}=W_t-\eta_tO_t,\qquad O_t\in\partial\nu(M_t),
\]
with $\beta\in[0,1)$ fixed, $O_t$ any measurable function of $(\xi_1,\dots,\xi_t)$ with values in $\partial\nu(M_t)$,
and $\eta_t=\eta_0(t+t_0)^{-\theta}$, $\eta_0,t_0>0$, $\frac12<\theta\le1$. Let
$N_\beta=(1-\beta)\sum_{j\ge0}\beta^j\xi_j$ be the stationary filtered noise. Then, almost surely, $\nabla F(W_t)$ converges
to a (possibly random) point
\[
G_\infty\in\arg\min\psi_{N_\beta}=-\Med_\nu(N_\beta),\qquad\text{and}\qquad W_t\to(\nabla F)^{-1}(G_\infty).
\]
No assumption is made on the law of the noise beyond a second moment: it may be discrete, lattice, or carried by a
proper subspace, and $\partial\nu(M_t)$ may be multivalued with positive probability at every step.}

\newcommand{\StmtCF}{%
Let $\nu$ be the $\ell_1$, the Euclidean or the nuclear norm with its direction map, and let Assumption~\ref{ass:N}
hold. Then for every $\epsilon>0$ there is $b_0$ such that for all $b\ge b_0$
\begin{equation}\label{eq:cf}
\arg\min\psi_{\xi_b}\ \subset\ B\Big(\frac{\gs}b,\frac\epsilon b\Big),\qquad\text{i.e.}\qquad
\Med_\nu(\xi_b)=-\frac{\gs}b+o\Big(\frac1b\Big),\qquad \gs=-\tfrac16\,J_\Sigma^{-1}\,T_\Sigma[K_3].
\end{equation}}

\newcommand{\StmtTwo}{%
Let $K,d\ge2$, $r=\min(K,d)$, $\bar d=\max(K,d)$, $\nu=\nnuc\cdot$ on $\R^{K\times d}$, and let $Z$ be a standard Gaussian
$K\times d$ matrix with singular values $s_1,\dots,s_r$. Put
\begin{equation}\label{eq:q}
m_p=\E\sum_is_i^p,\qquad q_{K,d}=\frac1{m_1}\,\E\sum_{i\ne j}\frac{s_is_j}{s_i+s_j},\qquad c_{K,d}=\frac{m_1}{Kd},
\end{equation}
\begin{equation}\label{eq:a12}
a_1=\frac{3\big[(K-1)(d-1)-2(K+d+1)\,q_{K,d}\big]}{(K-1)(d-1)(K+2)(d+2)},\qquad
a_2=\frac{6\big[(Kd+2)\,q_{K,d}+(K-1)(d-1)\big]}{(K-1)(d-1)(K+2)(d+2)} .
\end{equation}
\begin{enumerate}
\item[\textup{(a)}] \emph{(Reduction.)} For every symmetric $3$-tensor $K_3$ and $\sigma>0$,
$J_{\sigma^2I}=\sigma^{-1}c_{K,d}\,\Id$ and $T_{\sigma^2I}[K_3]=-\sigma^{-3}c_{K,d}\,(a_1\cA[K_3]+a_2\cB[K_3])$, where
$\cA[K_3]_{ij}=\sum_{k,l}K_3^{(ij),(kl),(kl)}$ and $\cB[K_3]_{ij}=\sum_{k,l}K_3^{(kj),(kl),(il)}$. Hence, for a noise law
with covariance $\sigma^2I$ and third cumulant $K_3=\E\zeta^{\otimes3}$,
\begin{equation}\label{eq:two}
\gs=\frac1{6\sigma^2}\Big(a_1\,\E\big[\zeta\nF\zeta^2\big]+a_2\,\E\big[\zeta\zeta^\top\zeta\big]\Big).
\end{equation}
The bias operator $K_3\mapsto\gs$ has rank $Kd$ on a space of dimension $\binom{Kd+2}3$.
\item[\textup{(b)}] \emph{(One number.)} $m_3=(K+d-q_{K,d})\,m_1$; $0<q_{K,d}<\frac{r-1}2$; $a_2>0$.
\item[\textup{(c)}] \emph{(Rank-one and full-rank directions.)}
$a_1+a_2=\dfrac{3(3+2q_{K,d})}{(K+2)(d+2)}$ and $r\,a_1+a_2=\dfrac{3(\bar d-1-2q_{K,d})}{(\bar d-1)(\bar d+2)}>0$.
\item[\textup{(d)}] \emph{(Negative fourth derivative.)} For every $D\ne0$,
$\nabla^4\psi_I(0)[D,D,D,D]=-c_{K,d}\big(a_1\nF D^4+a_2\tr(DD^\top)^2\big)<0$.
\item[\textup{(e)}] \emph{(Proportional limit.)} If $K,d\to\infty$ with $K/d\to\gamma\in(0,1]$, then $q_{K,d}/K\to\kappa(\gamma)$,
$d\,a_2\to6\kappa(\gamma)$ and $d^2a_1\to\frac3\gamma\big(1-2(1+\gamma)\kappa(\gamma)\big)$, where
$\kappa(\gamma)=\big[(1+\gamma)\mu_{1/2}(\gamma)-\mu_{3/2}(\gamma)\big]/\big(\gamma\,\mu_{1/2}(\gamma)\big)$ and
$\mu_s(\gamma)=\int x^s\,d\mathrm{MP}_\gamma(x)$ are the moments of the Marchenko--Pastur law; $\kappa(1)=\frac25$.
\end{enumerate}
For $K=1$ the two contractions coincide and $\gs=\frac{3}{d+2}\cdot\frac{\E[\zeta|\zeta|^2]}{6\sigma^2}$.}

\newcommand{\StmtAone}{%
For all $K,d\ge2$,
\[
q_{K,d}>\frac{(K-1)(d-1)}{2(K+d+1)},\qquad\text{equivalently}\qquad a_1<0<a_2,\qquad\text{equivalently}\qquad
a_1+a_2>\frac3{K+d+1}.
\]
With $W=ZZ^\top$, $D_p=\E\tr W^{p/2}$ and $r(p)=D_{p+2}/D_p$, each of the following is also equivalent:
\begin{enumerate}
\item[\textup{(i)}] $2\,r(1)+r(4)<3\,r(2)$: a three-point concavity of the moment ratio of a real Wishart matrix;
\item[\textup{(ii)}] $p\mapsto\E\sum_j s_j^pA_j/\E\sum_js_j^p$, with $A_j=\sum_{i\ne j}s_i^2$, is strictly convex at the
points $p=1,2,4$;
\item[\textup{(iii)}] for the size-biased mean spectral measure $\varpi(dx)\propto x\,\E\sum_i\delta_{s_i^2}(dx)$, of mean
$m=K+d+1$ and variance $v=Kd+K+d+3$: $\ \E_\varpi[x^{1/2}]/\E_\varpi[x^{-1/2}]<m-\frac v{2m}$.
\end{enumerate}
For $K=1$, and in \textup{(iii)} for every Gamma law, there is equality.}

\newcommand{\StmtDim}{%
For independent entries of variance $\sigma^2$ and third cumulant $\kappa_3$, $\cA=\cB=\kappa_3\ones\ones^\top$ and
\[
\gs=\frac{\kappa_3}{6\sigma^2}\,(a_1+a_2)\,\ones\ones^\top,\qquad
a_1+a_2=\frac{3\,(3+2q_{K,d})}{(K+2)(d+2)},\qquad
\frac{3}{K+d+1}<a_1+a_2<\frac{3}{\max(K,d)+2}\quad(K,d\ge2).
\]
Per entry, the first-order offset of a polar step is the scalar Cornish--Fisher offset $\kappa_3/(6\sigma^2b)$ of sign
descent divided by a factor between $\frac13(\max(K,d)+2)$ and $\frac13(K+d+1)$. In the proportional limit
$\max(K,d)(a_1+a_2)\to6\kappa(\gamma)$, equal to $\frac{12}5$ for square matrices.}

\newcommand{\StmtRankOne}{%
If $\zeta$ has rank one almost surely (per-sample gradients $\delta x^\top$ of a linear layer) and covariance
$\sigma^2I$, then $\zeta\zeta^\top\zeta=\nF\zeta^2\zeta$ and, for the nuclear norm,
$\gs=\frac{3+2q_{K,d}}{(K+2)(d+2)}\cdot\frac{\E[\zeta\nF\zeta^2]}{2\sigma^2}$, whereas for the Frobenius norm (normalized
steps) $\gs=\frac{1}{Kd+2}\cdot\frac{\E[\zeta\nF\zeta^2]}{2\sigma^2}$. Orthogonalization multiplies the first-order offset
of normalization by
\[
\varrho_{K,d}=\frac{(3+2q_{K,d})(Kd+2)}{(K+2)(d+2)},\qquad \frac{Kd+2}{K+d+1}<\varrho_{K,d}<\frac{Kd+2}{\max(K,d)+2}<\min(K,d).
\]}

\newcommand{\StmtCompare}{%
\textup{(i)} To first order a polar step sees the third cumulant only through $a_1\cA+a_2\cB$ ($Kd$ numbers), a
normalized step only through $\cA$, and sign descent only through the $Kd$ own third cumulants
$K_3^{(ij),(ij),(ij)}$.
\textup{(ii)} \emph{Sign-blind noise biases polar steps:} if $\zeta$ has covariance $\sigma^2I$ and all third cumulants zero
except $\E[\zeta_{11}\zeta_{12}^2]=\tau\ne0$ (and its permutations), sign descent has no first-order offset and
$\gs=\frac{(a_1+a_2)\tau}{6\sigma^2}e_1e_1^\top$.
\textup{(iii)} \emph{Polar-blind noise biases sign descent:} the tensor supported on $\{(11),(22)\}$ with
$K_3^{(11)^3}=1$, $K_3^{(11),(22),(22)}=-\frac{a_1+a_2}{a_1}$, $K_3^{(22)^3}=K_3^{(22),(11),(11)}=0$ has
$a_1\cA+a_2\cB=0$ and a nonzero own cumulant.
\textup{(iv)} \emph{One skewed factor:} if $\zeta=\epsilon M+\zeta'$ with $\epsilon$ a scalar of third cumulant
$\kappa_3(\epsilon)$, $M$ a fixed matrix, $\zeta'$ independent and symmetric, and total covariance $\sigma^2I$, then
$\gs=\frac{\kappa_3(\epsilon)}{6\sigma^2}(a_1\nF M^2M+a_2MM^\top M)$ and $\inner{\gs}{\kappa_3(\epsilon)M}>0$.}

\newcommand{\StmtBeff}{%
Let $\nu$ and the samples be as in Theorem~\ref{thm:cf}, and let $N_w=\sum_{k\ge0}w_k\xi_b^{(k)}$ for independent batch
means $\xi_b^{(k)}$ and fixed real weights with $\sum_k|w_k|<\infty$, $w\ne0$, $\rho_p=\sum_kw_k^p$. Then
\[
\arg\min\psi_{N_w}\subset B\Big(\frac{\rho_3}{\rho_2}\cdot\frac{\gs}b,\ o\Big(\frac1b\Big)\Big)\qquad(b\to\infty).
\]
For the exponential moving average $w_k=(1-\beta)\beta^k$ this is $\gs/\beff+o(1/b)$ with the \emph{skewness-effective
batch size}
\[
\beff=b\,\frac{1+\beta+\beta^2}{1-\beta^2},\qquad\text{not the variance-effective one}\quad b\,\frac{1+\beta}{1-\beta}.
\]}

\newcommand{\StmtState}{%
Let $\hat g_b(W)=\frac1b\sum_{i\le b}g(W;\omega_i)$ with $\nabla F(W)=\E g(W;\omega)$, let $\Sigma(W)$, $K_3(W)$ be the
covariance and third cumulant of $g(W;\omega)-\nabla F(W)$, and $h_b(W)=\E\,o(\hat g_b(W))$ the mean field. Assume, on a
bounded neighbourhood $\mathcal U$ of $W_\ast$: \textup{(i)} $F$ is $C^1$, twice differentiable at $W_\ast$,
$\nabla F(W_\ast)=0$, $H=\nabla^2F(W_\ast)$ is invertible and $|\nabla F(W)|\ge c|W-W_\ast|$; \textup{(ii)} $\Sigma(W)$,
$K_3(W)$ are bounded, continuous at $W_\ast$, and $\Sigma(W)\succeq\lambda I$; \textup{(iii)} the one-term expansion of
Lemma~\ref{lem:edge} holds uniformly in $W\in\mathcal U$. Let $\gs$ be computed from $\Sigma(W_\ast)$, $K_3(W_\ast)$.
Then for every $\epsilon>0$ there are $b_0$ and $\delta>0$ such that for $b\ge b_0$ every $W\in\mathcal U$ with
$|h_b(W)|\le\delta b^{-1/2}$, in particular every zero of the mean field in $\mathcal U$, satisfies
\[
\Big|W-W_\ast-\frac1b\,H^{-1}\gs\Big|\le\frac\epsilon b .
\]
Consequently $F(W)-F(W_\ast)=\frac1{2b^2}\inner{\gs}{H^{-1}\gs}+o(b^{-2})$ at such points when $H$ is symmetric positive
definite. If moreover $h_b$ is continuous on $\mathcal U$, it has a zero in $\mathcal U$ for $b$ large.}

\newcommand{\StmtLin}{%
Suppose $g(W;\omega)-\nabla F(W)=L_W\chi(\omega)$ on $\mathcal U$, where $\chi\in\R^p$ satisfies Assumption~\ref{ass:N}
and $W\mapsto L_W$ is continuous with $L_W:\R^p\to E$ onto. Then \textup{(ii)} and \textup{(iii)} of
Theorem~\ref{thm:state} hold on a neighbourhood of $W_\ast$.}

\newcommand{\StmtFS}{%
Let $F=\frac1n\sum_if_i$ on $\R^{K\times d}$, $(K,d)\ne(1,1)$, and at a point $W$ let $\zeta=\nabla f_I(W)-\nabla F(W)$ with
$I$ uniform on $\{1,\dots,n\}$. If the $n$ gradients $\nabla f_i(W)$ affinely span $\R^{K\times d}$ (so $n\ge Kd+1$),
then for minibatches of size $b$ drawn with replacement the nuclear medians of the batch noise are
$-\gs/b+O(b^{-1-\gamma}\log^{Kd}b)$ with $\gamma$ as in Theorem~\ref{thm:nuclear}, whatever the arithmetic of the
gradients. When $n=Kd+1$ the noise is a lattice law for \emph{every} data set: no weak Cram\'er or non-lattice
condition can hold, and the law holds nevertheless.}

\newcommand{\StmtLattice}{%
Let $E=\R^2$ with the Euclidean norm, and let $\zeta$ be supported on $a+\cL$ for a lattice $\cL$ of covolume $|\cL|$
generated by the support of $\zeta-\zeta'$. Let $\Lambda_b=ba+\cL$ and
\[
\mathfrak E_\Lambda(m)=\lim_{\tau\to\infty}\Big[|\cL|\sum_{\ell\in\Lambda}\frac{m+\ell}{|m+\ell|}e^{-\inner\ell{\Sigma^{-1}\ell}/2\tau^2}
-\int_{\R^2}\frac{m+y}{|m+y|}e^{-\inner y{\Sigma^{-1}y}/2\tau^2}dy\Big].
\]
Then $b\,G_b^\st=\gs-b^{-1/2}\,(2\pi\sqrt{\det\Sigma})^{-1}J_\Sigma^{-1}\mathfrak E_{\Lambda_b}(\gs)+o(b^{-1/2})$: the
second-order term of $G_b^\st$ is of exact order $b^{-3/2}$, and depends on $b$ only through the coset $\Lambda_b$. If so, the remainder of Theorem~\ref{thm:codim} is
attained for $c=2$.}
